% Figure: sound pressure level versus distance from a source in a room,
% decomposed into the direct field (6 dB per doubling) and the flat
% reverberant field, with the critical distance where they cross.
\begin{figure}[htbp]
\centering
\begin{tikzpicture}
\begin{axis}[
  width=12cm, height=7.5cm,
  xlabel={distance from source (m)}, ylabel={level above $L_W$ reference (dB)},
  xmode=log,
  xmin=0.2, xmax=20, ymin=-30, ymax=0,
  axis lines=left,
  domain=0.2:20, samples=300,
  legend pos=north east,
  legend cell align=left,
  every axis x label/.style={at={(current axis.right of origin)}, anchor=west},
  every axis y label/.style={at={(current axis.above origin)}, anchor=south},
  grid=major, grid style={enggray!20},
  tick label style={font=\scriptsize},
  label style={font=\small},
  legend style={font=\small},
]
% Direct field: L = 10 log10( Q / (4 pi r^2) ), Q = 2 (half-space).
\addplot[engblue, very thick, dashed] {10*log10( 2/(4*pi*x^2) )};
\addlegendentry{direct field ($1/r^2$)}
% Reverberant field: L = 10 log10( 4 / R ), room constant R = 50 m^2.
\addplot[engred, very thick, dashed] {10*log10(4/50)};
\addlegendentry{reverberant field (flat)}
% Total field: energy sum of the two.
\addplot[engblack, very thick] {10*log10( 2/(4*pi*x^2) + 4/50 )};
\addlegendentry{total field}
% critical distance where direct = reverberant, r_c = sqrt(Q R / (16 pi)).
\addplot[enggray, thin, dotted] coordinates {(1.41,-30) (1.41,0)};
\node[enggray, font=\scriptsize, anchor=south, rotate=90]
  at (axis cs:1.41,-22) {critical distance $r_c$};
\end{axis}
\end{tikzpicture}
\caption[Direct and reverberant field in a room]{Sound pressure level
versus distance from a source in a room, referenced to the source sound
power level, decomposed into the direct field and the reverberant field.
Close to the source the direct field dominates and the level falls six
decibels per doubling of distance, as in the open air. Far from the
source the reverberant field, built from the room's many reflections,
dominates and the level is nearly uniform, set by the room constant
$R = S\bar\alpha/(1-\bar\alpha)$. The two are equal at the critical
distance $r_c=\sqrt{Q R/(16\pi)}$, marked here for a directivity factor
$Q=2$ and $R=50$ m$^2$. The practical reading is direct: a listener or a
microphone closer than $r_c$ hears mostly the source and can be moved to
change the level, while beyond $r_c$ the level is fixed by the room, and
only added absorption (a larger $R$) lowers it. Sound-power surveys and
speech-intelligibility predictions both turn on which side of $r_c$ the
receiver sits.}
\label{fig:vol03ch07:room-field}
\end{figure}
